\Needspace{12\baselineskip}
\section{实验结果与分析}
\subsection{基准CNN与Dropout}
\textbf{无增强基准存在明显过拟合。}三种子末轮验证准确率为$83.15\pm0.29\%$，固定训练子集的评估准确率为$100.00\pm0.00\%$，泛化差距$16.85\pm0.29$个百分点。图\ref{fig:baseline}中，训练曲线逐渐饱和，验证曲线仍明显低于训练曲线。
\begin{figure}[htbp]\centering
\includegraphics[width=\linewidth]{baseline_curves.pdf}
\caption{基准、FC Dropout 0.5与BN的学习曲线（种子0）。实线为固定5000张训练子集的评估准确率，虚线为验证准确率。}\label{fig:baseline}
\end{figure}
\FloatBarrier
\begin{table}[htbp]\centering\small
\caption{Dropout对照：三种子均值$\pm$标准差，验证损失为三种子均值。}\label{tab:dropout}
\ReportTable{tables/dropout}
\end{table}

表\ref{tab:dropout}中，FC Dropout由0增至0.7时，验证准确率逐步提高；0.7配置比基准高2.61个百分点。随机丢弃限制了隐藏特征间的共同适应，缩小了训练与验证差距。卷积0.1+FC0.5为$85.23\pm0.18\%$，与单用FC0.5的表现接近。卷积0.2组合的验证损失降至0.598，准确率为85.09\%，说明预测概率的改善与类别判断的改善并不同步。
\begin{figure}[htbp]\centering
\includegraphics[width=\linewidth]{dropout_gap.pdf}
\caption{Dropout的验证表现。浅色点为三个种子，菱形及误差条为均值$\pm$标准差。}\label{fig:dropout}
\end{figure}

\FloatBarrier
\subsection{归一化及其与训练配方的相互作用}
表\ref{tab:norm}在固定学习率0.02、batch128、30轮条件下比较归一化。BN验证均值86.77\%，比无归一化高3.62个百分点；GN、LN、IN分别为85.95\%、85.35\%、84.52\%。归一化使中间激活的尺度更稳定，改善了当前配方的收敛与泛化；其中BN的验证均值最高。
\begin{table}[htbp]\centering\small
\caption{归一化对照：三种子均值$\pm$标准差。}\label{tab:norm}
\ReportTable{tables/norm_report}
\end{table}
\begin{figure}[htbp]\centering
\includegraphics[width=\linewidth]{norm_curves.pdf}
\caption{各归一化方法的验证收敛曲线（种子0）。}\label{fig:norm}
\end{figure}

将batch从128改为16，并将学习率线性降至0.0025后，三种归一化的验证变化均在0.3个百分点以内（表\ref{tab:batch}）。BN在这组小批量配方下仍保持相近表现；GN与LN的统计量按样本计算，不依赖批内样本数。
\begin{table}[htbp]\centering\small
\caption{批大小配方比较，均为种子0；batch16同时采用八分之一学习率。}\label{tab:batch}
\begin{tabular}{lccc}
\toprule
归一化 & Batch 128 (\%) & Batch 16 (\%) & 变化 (pp) \\
\midrule
BatchNorm & 86.68 & 86.86 & +0.18 \\
GroupNorm & 85.76 & 86.02 & +0.26 \\
LayerNorm & 85.82 & 85.66 & -0.16 \\
\bottomrule
\end{tabular}

\end{table}

学习率对照见表\ref{tab:lr}及图\ref{fig:lr}。峰值0.1时none/BN/GN分别为85.22\%、89.00\%、87.28\%；峰值0.4配方下，无归一化网络在第1轮预热阶段出现NaN，而BN与GN仍得到88.50\%与84.82\%。归一化拓宽了本配方的可训练学习率范围；GN在0.4下准确率回落，表明步长增大后优化效果也会下降。
\begin{table}[htbp]\centering\small
\caption{不同峰值学习率的末轮验证准确率（\%），种子0；0.02列复用基础对照。}\label{tab:lr}
\begin{tabular}{lcccc}
\toprule
归一化 & lr=0.005 & lr=0.02 & lr=0.1 & lr=0.4 \\
\midrule
无 & 79.16 & 83.46 & 85.22 & 发散 \\
BatchNorm & 81.98 & 86.68 & 89.00 & 88.50 \\
GroupNorm & 81.72 & 85.76 & 87.28 & 84.82 \\
\bottomrule
\end{tabular}

\end{table}
\begin{figure}[htbp]\centering
\includegraphics[width=\linewidth]{norm_lr.pdf}
\caption{学习率敏感性（种子0）。无归一化的0.4任务在第1轮出现NaN。}\label{fig:lr}
\end{figure}

表\ref{tab:combo}中，BN加FC Dropout0.5由$86.77\pm0.16\%$提高至$88.72\pm0.45\%$，再加卷积Dropout0.1则为$87.83\pm0.08\%$。BN与FC Dropout在本组实验中形成互补，而进一步增加卷积通道丢弃降低了验证均值。卷积层和分类器上的正则化强度需要匹配。
\begin{table}[htbp]\centering\small
\caption{Dropout与BN组合的末轮验证准确率（\%），三种子均值$\pm$样本标准差。}\label{tab:combo}
\begin{tabular}{lccc}
\toprule
归一化 & 无 Dropout & FC 0.5 & Conv 0.1 + FC 0.5 \\
\midrule
无 & 83.15 $\pm$ 0.29 & 85.07 $\pm$ 0.19 & 85.23 $\pm$ 0.18 \\
BatchNorm & 86.77 $\pm$ 0.16 & 88.72 $\pm$ 0.45 & 87.83 $\pm$ 0.08 \\
\bottomrule
\end{tabular}

\end{table}

\FloatBarrier
\Needspace{7\baselineskip}
\subsection{搜索、五折选型与最终测试}
24组搜索结果完整列于附录表\ref{tab:search}。前4名为cfg00、cfg02、cfg08、cfg18，15轮验证准确率分别为89.42\%、88.64\%、88.42\%、88.34\%。cfg01采用GN、学习率0.1578和FC Dropout0.5，损失有限但末轮验证仅10.00\%，表现为优化塌缩。
\begin{table}[htbp]\centering\footnotesize
\caption{四候选的五折验证准确率（\%）：30轮、种子0，报告折间标准差。}\label{tab:cv}
\begin{tabular}{lcccccc}
\toprule
配置 & 折1 & 折2 & 折3 & 折4 & 折5 & 均值 $\pm$ STD \\
\midrule
cfg00 & 88.73 & 88.53 & 88.50 & 88.22 & 88.57 & 88.51 $\pm$ 0.18 \\
cfg08 & 88.42 & 88.39 & 88.39 & 88.06 & 88.68 & 88.39 $\pm$ 0.22 \\
cfg02 & 87.95 & 87.78 & 88.55 & 88.04 & 88.33 & 88.13 $\pm$ 0.31 \\
cfg18 & 87.58 & 87.53 & 87.56 & 87.42 & 87.76 & 87.57 $\pm$ 0.12 \\
\bottomrule
\end{tabular}

\end{table}

表\ref{tab:cv}中cfg00均值$88.51\pm0.18\%$，按预定规则被选中；cfg08为$88.39\pm0.22\%$，均值差为0.12个百分点。cfg00的精确配置为BN、batch128、学习率0.12089、weight decay=0.0009851、FC Dropout0.3。
\begin{table}[htbp]\centering\small
\caption{最终测试结果。折模型行报告五模型均值$\pm$标准差；A1为历史参考。}\label{tab:final}
\ReportTable{tables/final}
\end{table}

在全部50000张训练图上重训后，单模型测试89.25\%、NLL0.405；五折模型概率集成测试90.20\%、NLL0.319，较各折单模型的均值88.28\%更高。概率平均综合了不同折模型的判断，降低单个模型预测误差的影响，同时增加推理成本。
\FloatBarrier
